% Vietnamese translation for OI Public Library
% Source-File: IOI中国国家候选队论文/2008/Day2/9.苏煜《对块状链表的一点研究》/对块状链表的一点研究.doc
\documentclass[11pt]{article}
\usepackage{oipl-vn}
\usepackage{tikz}
\usetikzlibrary{positioning}

\newcommand{\Oh}{\mathcal{O}}

\title{Một vài nghiên cứu về danh sách liên kết theo khối}
\author{Su Yu\\Trường Trung học trực thuộc Đại học Sơn Tây}
\date{Trại đông Tin học toàn quốc, 2008}

\begin{document}
\maketitle

\origtitle{Một vài nghiên cứu về danh sách liên kết theo khối}
\sourcepath{IOI China candidate team papers / 2008 / Day 2 / Su Yu / block linked list study.doc}

\renewcommand{\abstractname}{Tóm tắt}
\begin{abstract}
Bài viết giới thiệu khái niệm danh sách liên kết theo khối, cách mở rộng cấu
trúc này, phân tích hiệu năng và bàn về mặt lợi hại khi dùng nó trong các kỳ
thi tin học. Phần cuối nêu thêm một vài ứng dụng thực tế của tư tưởng chia
khối.
\end{abstract}

\noindent\textbf{Từ khóa:} danh sách liên kết theo khối; kích thước khối;
hiệu năng; mở rộng cấu trúc dữ liệu; mô phỏng; lấy điểm từng phần.

\section{Danh sách liên kết theo khối là gì?}

Ta bắt đầu từ một bài toán để thấy cấu trúc này xuất hiện tự nhiên.

\subsection{NOI 2003: Editor}

Bài toán yêu cầu xây một trình soạn thảo văn bản rỗng ban đầu. Văn bản là một
dãy gồm các ký tự ASCII nhìn thấy được, con trỏ có thể nằm ở đầu, cuối hoặc
giữa hai ký tự. Chương trình phải xử lý sáu thao tác sau.

\begin{center}
\begin{tabular}{lll}
\textbf{Thao tác} & \textbf{Dạng trong input} & \textbf{Ý nghĩa}\\
\hline
\texttt{MOVE(k)} & \texttt{Move k} & đưa con trỏ tới sau ký tự thứ \(k\)\\
\texttt{INSERT(n,s)} & \texttt{Insert n} rồi chuỗi \(s\) & chèn \(n\) ký tự tại con trỏ\\
\texttt{DELETE(n)} & \texttt{Delete n} & xóa \(n\) ký tự sau con trỏ\\
\texttt{GET(n)} & \texttt{Get n} & xuất \(n\) ký tự sau con trỏ\\
\texttt{PREV()} & \texttt{Prev} & lùi con trỏ một ký tự\\
\texttt{NEXT()} & \texttt{Next} & tiến con trỏ một ký tự
\end{tabular}
\end{center}

Ví dụ, từ trình soạn thảo rỗng, lần lượt thực hiện
\texttt{INSERT(13, "Balanced tree")}, \texttt{MOVE(2)}, \texttt{DELETE(5)},
\texttt{NEXT()}, \texttt{INSERT(7, " editor")}, \texttt{MOVE(0)},
\texttt{GET(16)} thì kết quả là \texttt{"Bad editor tree"}.

Giới hạn chính của bài là: số thao tác \texttt{MOVE} không quá \(50000\), tổng
số thao tác \texttt{INSERT} và \texttt{DELETE} không quá \(4000\), tổng số
\texttt{PREV} và \texttt{NEXT} không quá \(200000\). Tổng số ký tự được chèn
không vượt quá \(2\) MB, còn output đúng không vượt quá \(3\) MB. Mọi thao tác
đều hợp lệ.

Thực chất các lệnh chỉ thuộc hai nhóm: định vị, và thêm hoặc xóa. Hai cách
lưu trữ tuần tự quen thuộc có ưu nhược điểm đối lập:
\[
\begin{array}{c|cc}
 & \text{mảng} & \text{danh sách liên kết}\\
\hline
\text{định vị} & \Oh(1) & \Oh(N)\\
\text{thêm hoặc xóa} & \Oh(N) & \Oh(1)
\end{array}
\]
Vì tồn tại thao tác \(\Oh(N)\), dùng riêng một trong hai cấu trúc đều không
đủ tốt. Nếu kết hợp chúng, tức nhìn toàn cục như một danh sách liên kết nhưng
mỗi nút chứa một mảng nhỏ có kích thước thích hợp, chẳng hạn \(1000\) hoặc
\(1500\), ta nhận được một cấu trúc cân bằng hơn: đó là
\term{danh sách liên kết theo khối}.

Với một biến nguyên ghi vị trí hiện tại, ta cần cấu trúc hỗ trợ ba thao tác:
\begin{enumerate}
  \item chèn một đoạn thông tin tại vị trí chỉ định;
  \item xóa một đoạn thông tin từ vị trí chỉ định;
  \item lấy một đoạn thông tin từ vị trí chỉ định.
\end{enumerate}

Cài đặt cơ bản gồm hai thao tác nội bộ:
\begin{itemize}
  \item \term{định vị}: đi từ khối đầu tiên về sau cho tới khi tìm được khối
  chứa vị trí cần tìm, cùng vị trí tương đối trong khối đó;
  \item \term{tách khối}: tách một khối tại một vị trí thành hai khối.
\end{itemize}
Từ đó, thao tác chèn tìm vị trí, tách khối, rồi thêm các khối mới cho tới khi
chèn xong. Thao tác xóa tìm đầu đoạn, tách khối đầu và khối cuối, rồi bỏ toàn
bộ các nút nằm giữa. Thao tác lấy đoạn cũng định vị điểm bắt đầu, sau đó quét
qua các khối cho tới khi đủ số ký tự.

Hai thao tác chính có thể hình dung như sau. Với thao tác chèn, ta tách khối
chứa vị trí bắt đầu rồi móc các khối mới vào giữa hai nửa. Với thao tác xóa,
ta tách hai biên của đoạn cần xóa, bỏ các khối nằm trọn trong đoạn đó, rồi nối
phần trước với phần sau.

\begin{figure}[ht]
\centering
\begin{tikzpicture}[
  block/.style={draw, minimum width=1.25cm, minimum height=0.62cm, align=center},
  small/.style={draw, minimum width=0.85cm, minimum height=0.62cm, align=center},
  arrow/.style={->, thick},
  node distance=0.08cm
]
\node[block] (a) {A};
\node[small, right=0.35cm of a] (b1) {B\(_L\)};
\node[small, right=0.08cm of b1] (b2) {B\(_R\)};
\node[block, right=0.35cm of b2] (c) {C};
\node[block, below=0.8cm of b1] (x) {mới};
\node[block, right=0.08cm of x] (y) {mới};
\draw[arrow] (a) -- (b1);
\draw[arrow] (b1) -- (x);
\draw[arrow] (x) -- (y);
\draw[arrow] (y) -- (b2);
\draw[arrow] (b2) -- (c);
\draw[dashed] (b1.north east) -- (b1.south east);
\node[above=0.18cm of b1] {\small tách tại vị trí chèn};
\end{tikzpicture}
\caption{Chèn trong danh sách khối: tách khối hiện tại và nối các khối mới vào giữa.}
\end{figure}

\begin{figure}[ht]
\centering
\begin{tikzpicture}[
  block/.style={draw, minimum width=1.15cm, minimum height=0.62cm, align=center},
  del/.style={draw, dashed, fill=gray!10, minimum width=1.15cm, minimum height=0.62cm, align=center},
  arrow/.style={->, thick}
]
\node[block] (a) {A};
\node[block, right=0.35cm of a] (l) {L};
\node[del, right=0.35cm of l] (m1) {xóa};
\node[del, right=0.08cm of m1] (m2) {xóa};
\node[block, right=0.35cm of m2] (r) {R};
\node[block, right=0.35cm of r] (c) {C};
\draw[arrow] (a) -- (l);
\draw[arrow] (l) to[bend right=18] node[below] {\small nối lại} (r);
\draw[arrow] (r) -- (c);
\draw[<->, thick] (m1.north west) -- node[above] {\small đoạn bị xóa} (m2.north east);
\end{tikzpicture}
\caption{Xóa trong danh sách khối: sau khi tách hai biên, bỏ các khối ở giữa và nối lại.}
\end{figure}

Một vấn đề còn lại là việc tách khối quá thường xuyên có thể tạo ra nhiều
khối liên tiếp chứa rất ít dữ liệu, làm giảm hiệu quả. Vì vậy sau mỗi thao
tác ta có thể gộp các khối liên tiếp quá nhỏ. Như vậy ta đã có một danh sách
liên kết theo khối cơ bản đủ đáp ứng yêu cầu của bài.

\section{Ứng dụng đơn giản}

Danh sách liên kết theo khối có thể xem là một cách tối ưu hóa mô phỏng. Nó
ghép hai cấu trúc không hoàn hảo thành một cấu trúc thực dụng hơn. Nhờ đặc
điểm đó, với một số bài ta có thể mô phỏng trực tiếp một cách khá ``cưỡng
bức''.

\subsection{NEERC 2003 Northern Subregion: Key Insertion}

Có \(N\) binh sĩ đang tập đội hình và \(M\) vị trí từ trái sang phải, với
\(1\le N,M\le 131072\). Mỗi lệnh có dạng \(\texttt{Goto}(L,S)\). Nếu vị trí
\(L\) trống thì binh sĩ \(S\) đứng vào đó. Nếu vị trí \(L\) đã có binh sĩ
\(K\), binh sĩ \(S\) đứng vào \(L\), rồi tiếp tục thực hiện
\(\texttt{Goto}(L+1,K)\). Mỗi binh sĩ nhận đúng một lệnh. Cần tìm trạng thái
cuối cùng của đội hình; trong quá trình đó chỉ số vị trí có thể vượt quá
\(M\), nên trước hết phải xác định độ dài cuối cùng. Vị trí trống được in
bằng \(0\).

Một cách hơi thô là nhìn thao tác chèn như sau: đẩy cả đoạn liên tiếp bắt đầu
tại vị trí hiện tại sang phải một ô, rồi đặt phần tử mới vào chỗ đó. Tương
đương, ta xóa một ô trống ngay sau đoạn liên tiếp ấy và chèn một phần tử mới
ở vị trí hiện tại. Việc duy trì ô trống đầu tiên sau một vị trí là ứng dụng
kinh điển của DSU, còn chèn xóa trên dãy lớn là điểm mạnh của danh sách liên
kết theo khối. Do đó có thể ghép DSU với danh sách khối để giải bài.

Lời giải chính thống của bài tinh tế hơn. Nó dùng DSU duy trì các khối vị trí
không kề nhau, và dùng danh sách liên kết lưu các phần tử trong mỗi khối. Khi
binh sĩ \(A\) được chèn vào một vị trí đã thuộc khối \(B\), ta đặt \(A\) vào
đầu danh sách của khối \(B\). Nếu thao tác chèn làm một hoặc nhiều khối nối
lại với nhau, danh sách của khối phía sau được nối vào trước danh sách của
khối phía trước. Sau khi xử lý xong, duyệt từ khối cuối cùng về trước; với
mỗi phần tử \(x\), chèn nó vào ô trống thứ \(L[x]\). Như vậy thu được dãy
cuối cùng.

Chỉ riêng mô tả cách này đã khá phức tạp, chưa nói tới việc nghĩ ra trong
phòng thi. Thực tế năm đó không đội nào giải được bài. Trong khi đó, với danh
sách liên kết theo khối, cách làm mô phỏng gần như không cần nhiều suy nghĩ;
có thể xem như một cách lấy điểm bằng mô phỏng mạnh hơn.

\section{Phân tích hiệu năng: chọn kích thước khối}

Trong trường hợp lý tưởng, giả sử kích thước mảng nhỏ trong mỗi nút là \(x\).
Khi đó có khoảng \(n/x\) khối. Thao tác định vị tốn \(\Oh(n/x)\), còn phần
ảnh hưởng phụ khi thêm hoặc xóa trong một khối tốn \(\Oh(x)\). Hai chi phí
cân bằng nhất khi
\[
  x \approx \sqrt n,
\]
và độ phức tạp của mỗi thao tác vào khoảng \(\Oh(\sqrt n)\). Trong thực tế
\(n\) thay đổi liên tục, nhưng chọn \(x\) theo cỡ \(\sqrt N\), với \(N\) là
cận trên của số phần tử, vẫn giữ được hiệu năng xấp xỉ \(\Oh(\sqrt N)\) trong
trường hợp xấu. Nói cách khác, dùng danh sách khối là theo đuổi sự cân bằng.

Bản gốc có hai biểu đồ thời gian: một cho \(10\) test của \textit{NOI 2003
Editor}, một cho \(38\) test của \textit{Key Insertion}. Các hình không được
chép lại ở đây; nội dung quan trọng là quan sát của tác giả: kích thước khối
lý thuyết không phải luôn chạy nhanh nhất, còn kích thước lớn hơn giá trị lý
thuyết thường cho tốc độ tốt hơn.

Nguyên nhân đến từ cách cài đặt. Thao tác chèn và xóa đều phụ thuộc vào tách
khối, nên việc tách diễn ra thường xuyên và đa số khối không đầy. Kích thước
lưu trữ thực tế trong mỗi khối nhỏ hơn dung lượng tối đa. Vì vậy để kích
thước thực tế gần \(\sqrt n\), dung lượng tối đa nên lớn hơn \(\sqrt n\).
Mặt khác, vì hai khối kề nhau đều dưới nửa dung lượng sẽ được gộp, dung lượng
tối đa cũng không cần lớn hơn \(2\sqrt n\).

Với \textit{Editor}, tác giả nhận thấy dung lượng khối khoảng \(6000\) mới
nhanh nhất, lệch khá xa phân tích lý thuyết. Lý do là nhiều test chứa rất
nhiều nhóm lệnh như:
\[
\begin{array}{lll}
\texttt{MOVE x}, & \texttt{PREV}, & \texttt{NEXT},\\
\texttt{GET 1}, & \texttt{GET 1}, & \texttt{GET 1}.
\end{array}
\]
Những nhóm này không thay đổi cấu trúc khối; với \texttt{GET 1}, chi phí lấy
ký tự gần như không phụ thuộc kích thước khối. Nhưng trước khi lấy ký tự vẫn
cần định vị, nên chúng thực chất làm tăng tỷ trọng thao tác định vị. Khối càng
lớn thì số khối càng ít, định vị càng nhanh. Ngoài ra đề giới hạn tổng số
\texttt{INSERT} và \texttt{DELETE} chỉ \(4000\), nên kích thước khối không
ảnh hưởng mạnh bằng chi phí định vị.

Tác giả cũng viết một phiên bản dùng hai biến để ghi khối hiện tại và vị trí
tương đối trong khối. Vì cài đặt dùng danh sách một chiều, thao tác
\texttt{PREV} vẫn khó tránh khỏi định vị lại. Phiên bản này nhanh hơn, kích
thước tối ưu giảm còn khoảng \(5000\), và khối càng lớn thì mức cải thiện càng
nhỏ. Điều đó cho thấy một phần lớn lợi thế của khối lớn trong phiên bản đầu
đến từ chi phí định vị. Nếu dùng danh sách hai chiều và xây một đối tượng gần
giống iterator, khối nhỏ có thể nhanh hơn nữa, nhưng mã sẽ dài hơn. Tác giả
không đưa iterator vào cấu trúc cơ bản vì nó không phải lúc nào cũng hữu ích,
chẳng hạn \textit{Key Insertion} không cần, còn việc duy trì iterator làm tăng
độ dài mã và khả năng lỗi.

Với \textit{Key Insertion}, thời gian khi kích thước khối từ \(400\) tới
\(1200\) gần như nhau, tốt nhất quanh \(800\), tức nằm giữa \(\sqrt n\) và
\(2\sqrt n\). Tóm lại, hiệu năng danh sách khối phụ thuộc dữ liệu thực tế,
nhưng thường nên chọn dung lượng khối lớn hơn \(\sqrt n\) một chút để lượng
dữ liệu thực trong mỗi khối gần \(\sqrt n\). Đây vẫn là nguyên tắc cân bằng.

Do tách khối thường xuyên kéo theo cấp phát và giải phóng bộ nhớ, toàn bộ các
chương trình của tác giả dùng mảng tĩnh và tự duy trì danh sách nút rảnh.

\section{Mở rộng danh sách liên kết theo khối}

Xét một bài phức tạp hơn.

\subsection{NOI 2005: Sequence}

Cần duy trì một dãy số và hỗ trợ sáu thao tác:
\begin{center}
\begin{tabular}{lll}
\textbf{Số} & \textbf{Thao tác} & \textbf{Ý nghĩa}\\
\hline
1 & \texttt{INSERT posi tot c1 ... ctot} & chèn \(tot\) số sau vị trí \(posi\)\\
2 & \texttt{DELETE posi tot} & xóa \(tot\) số từ vị trí \(posi\)\\
3 & \texttt{MAKE-SAME posi tot c} & gán đoạn dài \(tot\) thành \(c\)\\
4 & \texttt{REVERSE posi tot} & đảo đoạn dài \(tot\)\\
5 & \texttt{GET-SUM posi tot} & hỏi tổng đoạn dài \(tot\)\\
6 & \texttt{MAX-SUM} & hỏi tổng lớn nhất của một đoạn con liên tiếp
\end{tabular}
\end{center}

Tại mọi thời điểm dãy có ít nhất một số, input luôn hợp lệ. Với \(50\%\) dữ
liệu, dãy có nhiều nhất \(30000\) số; với toàn bộ dữ liệu, dãy có nhiều nhất
\(500000\) số. Mỗi số nằm trong \([-1000,1000]\), số thao tác \(M\le 20000\),
tổng số phần tử được chèn không quá \(4000000\), input không quá \(20\) MB.

Mô phỏng trực tiếp thì dễ nhưng không đủ nhanh, và bài cũng không gợi ra một
thuật toán hoàn toàn khác. Đây là lúc danh sách liên kết theo khối lại hữu
dụng. So với bài trước, ta cần thêm các thao tác:
\begin{enumerate}
  \item đảo một đoạn con;
  \item gán cả một đoạn con thành cùng một giá trị;
  \item tính tổng một đoạn liên tiếp;
  \item tính tổng đoạn con liên tiếp lớn nhất của toàn dãy.
\end{enumerate}

Ưu điểm của danh sách khối là xử lý nguyên khối. Vì vậy mỗi khối cần lưu thêm
thông tin để các thao tác trên đoạn có thể bỏ qua các khối nằm trọn trong
đoạn. Để hỏi tổng đoạn con lớn nhất, mỗi khối lưu tổng của khối, tổng đoạn con
lớn nhất bên trong khối, và tổng lớn nhất của đoạn con chứa đầu trái hoặc đầu
phải của khối. Để gán nhanh cả khối thành một giá trị, dùng một cờ đánh dấu
khối đồng nhất và một biến lưu giá trị đó. Tương tự, mỗi khối có thêm cờ đảo
để biểu diễn trạng thái đang bị lật.

Lúc này mới thấy lợi ích của cách cài đặt trước đó: chèn và xóa không cần xét
quá nhiều thông tin mới, còn \texttt{REVERSE} và \texttt{MAKE-SAME} có thể
dùng cùng khuôn xử lý. Việc cập nhật thông tin tập trung chủ yếu vào thao tác
tách khối, nên chỉ cần viết phần đó thật cẩn thận.

Một mẹo nhỏ là lưu hai giá trị biên trong mảng \texttt{sidemax[2]}:
\texttt{sidemax[0]} là tổng lớn nhất của đoạn con chứa đầu trái của khối khi
chưa xét trạng thái đảo, còn \texttt{sidemax[1]} tương tự cho đầu phải. Khi
đó \texttt{sidemax[reversed]} chính là giá trị thực tế chứa đầu trái, và
\texttt{sidemax[!reversed]} là giá trị thực tế chứa đầu phải.

\section{Các ứng dụng phức tạp hơn}

\subsection{CERC 2007: Sort}

Trong một xưởng có \(N\) linh kiện xếp thành hàng, \(1\le N\le 100000\), biết
chiều cao của từng linh kiện. Cần sắp xếp chúng tăng dần, nhưng chỉ được dùng
thao tác sau: tìm vị trí \(P_1\) của linh kiện thấp nhất và đảo đoạn
\([1,P_1]\), rồi tìm vị trí \(P_2\) của linh kiện thấp thứ hai và đảo
\([2,P_2]\), cứ tiếp tục như vậy. Chương trình phải in ra \(P_1,P_2,P_3,\ldots\).
Nếu nhiều linh kiện có cùng chiều cao, xử lý linh kiện đứng trước trong dãy
ban đầu trước.

Thoạt nhìn khó tìm ngay một thuật toán đặc biệt bằng cây cân bằng, cây đoạn
hay heap. Vậy cứ mô phỏng. Ta vẫn dùng danh sách khối: vì cần đảo đoạn nên
thêm cờ đảo; vì mỗi bước cần tìm phần tử nhỏ nhất của dãy hiện tại nên mỗi
khối lưu thêm giá trị nhỏ nhất của khối. Ở mỗi bước, tìm khối chứa giá trị
nhỏ nhất cần xử lý, tìm vị trí cụ thể trong khối, đảo đoạn tiền tố, xóa phần
tử đầu, rồi lặp lại.

Vấn đề duy nhất là các giá trị bằng nhau và yêu cầu ưu tiên phần tử xuất hiện
sớm hơn trong dãy gốc. Có thể khử vấn đề này bằng cách thay giá trị của mỗi
phần tử bằng thứ tự xử lý của nó trong dãy ban đầu. Khi đó không còn trùng
giá trị, phần còn lại là cài đặt danh sách khối.

\subsection{NOI 2007: Necklace}

Một hệ thống cho phép khách hàng tự thiết kế dây chuyền. Dây chuyền có
\(N\) hạt, \(1\le N\le 500000\), mỗi hạt có màu thuộc \(1,2,\ldots,c\). Dây
được cố định trên một tấm phẳng; vị trí đánh dấu là \(1\), các vị trí còn lại
được đánh số theo chiều kim đồng hồ.

Hệ thống cần hỗ trợ:
\begin{itemize}
  \item \texttt{R k}: quay dây chuyền theo chiều kim đồng hồ \(k\) vị trí;
  \item \texttt{F}: lật tấm phẳng theo trục đối xứng cho sẵn;
  \item \texttt{S i j}: đổi chỗ hai hạt ở vị trí \(i,j\);
  \item \texttt{P i j x}: tô đoạn từ \(i\) tới \(j\) theo chiều kim đồng hồ
  thành màu \(x\);
  \item \texttt{C}: hỏi toàn bộ dây chuyền hiện có bao nhiêu đoạn màu liên
  tiếp;
  \item \texttt{CS i j}: hỏi đoạn từ \(i\) tới \(j\) theo chiều kim đồng hồ
  gồm bao nhiêu đoạn màu liên tiếp.
\end{itemize}

Thông tin cần duy trì lúc này là số đoạn màu. Để thuận tiện, mỗi khối còn lưu
màu hai đầu. Lật và quay có thể quy về thao tác đảo trước đó; truy vấn
\texttt{C} có thể dùng \texttt{CS}. Vì vậy cần thêm hai thao tác chính:
\texttt{Swap} và \texttt{CountSegment}. \texttt{Swap} chỉ cần tìm hai phần tử
cần đổi; nếu chúng khác nhau thì đổi rồi cập nhật thông tin của các khối liên
quan. \texttt{CountSegment} giống truy vấn tổng đoạn: khối đầu và cuối xử lý
riêng, các khối nằm giữa dùng thông tin đã lưu.

\section{Độ phức tạp cài đặt và tốc độ}

Với \textit{CERC 2007 Sort}, mã của tác giả dài \(3.44\) KB nhưng không nhanh.
Do không biết thời hạn của kỳ thi, tác giả so sánh với chương trình chuẩn dài
\(4.94\) KB, chứa nhiều chú thích và khoảng trắng: chương trình chuẩn chạy
\(0.61\) giây, còn chương trình danh sách khối chạy \(3.63\) giây.

Với \textit{NOI 2007 Necklace}, mã danh sách khối dài \(6.34\) KB và chỉ qua
\(7\) test. Nếu nhận ra thao tác lật và quay là những phép biến đổi tọa độ
``đánh lừa'', thêm tối ưu có thể qua \(9\) test. Tuy nhiên điều này cũng cho
thấy với người bình thường, viết danh sách khối đủ nhanh không dễ. Nếu đã nghĩ
đến biến đổi tọa độ để xử lý lật và quay, ta cũng có thể nghĩ tới cây đoạn,
đơn giản và nhanh hơn. Chênh lệch giữa \(\Oh(\sqrt n)\) và \(\Oh(\log n)\) là
hàng chục lần, hằng số của danh sách khối lại lớn hơn cây đoạn. Mã Pascal cây
đoạn của Yu Linyun chỉ dài \(4.12\) KB và nhanh hơn rất nhiều.

\section{Kết luận}

Ban đầu tác giả muốn giới thiệu danh sách liên kết theo khối như một cách lấy
điểm bằng mô phỏng, vì mô phỏng giúp tiết kiệm thời gian suy nghĩ. Nhưng cuối
cùng hiệu quả lấy điểm không tốt như mong đợi. Trước hết, mã danh sách khối
thường khá dài và dễ sai. Các bài hiện nay hay kết hợp nhiều cấu trúc dữ liệu,
như \textit{Key Insertion}; yêu cầu biến đổi trước khi mô phỏng, như
\textit{Sort}; hoặc cần cấu trúc có nhiều chức năng, như \textit{Sequence} và
\textit{Necklace}. Những tình huống đó làm việc đảm bảo đúng rất khó.

Thứ hai, ngay cả khi viết nhanh và đúng, hiệu năng vẫn không chắc chắn. Chọn
kích thước khối không tốt có thể vượt thời gian; đôi khi chọn gần tối ưu vẫn
có thể không đủ, vì trên dữ liệu lớn \(\Oh(\sqrt n)\) và \(\Oh(\log n)\), thậm
chí \(\Oh(1)\) trong một số bài, cách nhau rất xa.

Dẫu vậy, danh sách khối không hề vô dụng. Nó kết hợp mảng và danh sách liên
kết, hai cấu trúc cơ bản đều có khuyết điểm, rồi xử lý theo khối để cân bằng
hai loại thao tác. Tư tưởng ghép các phần tử cơ bản, thậm chí các phần tử cơ
bản không thật nhanh, để tạo thành kết quả hiệu quả hơn vẫn rất đáng học.

Hơn nữa, danh sách khối dễ mở rộng. Với thao tác trên dãy như gán đồng loạt,
đảo đoạn, hỏi tổng, duy trì giá trị nhỏ nhất, v.v., ta thường có thể đạt
\(\Oh(\sqrt n)\). Nếu duy trì toàn bộ danh sách khối theo thứ tự, nó còn có
thể mô phỏng một số thao tác của cây cân bằng, vì cây cân bằng cũng có thể
được nhìn như một cách duy trì dãy.

Do mỗi khối chỉ lưu một lượng thông tin phụ, mức sử dụng bộ nhớ khá tốt. Cùng
là mô phỏng, splay phải lưu toàn bộ thông tin phụ ở từng nút; tốc độ nhanh hơn
nhưng tốn bộ nhớ hơn.

Trong đời sống, tư tưởng danh sách khối cũng xuất hiện thường xuyên. Hệ thống
tập tin FAT truyền thống chia các sector đĩa thành cluster, rồi dùng bảng FAT
(\textit{File Allocation Table}) ghi trạng thái của từng cluster: có hỏng
không, có đang được dùng không, và nếu được dùng thì cluster kế tiếp là gì.
Cách tổ chức đó rất gần với danh sách liên kết theo khối.

Vì danh sách khối tiết kiệm không gian và cấu trúc khối dễ kết hợp với bộ
đệm, Vim cũng dùng cấu trúc tương tự cho lưu trữ trong bộ nhớ và bộ đệm trên
đĩa. Nếu dùng splay để viết trình soạn thảo, cấu trúc sẽ phức tạp và trừu
tượng hơn nhiều; cũng khó tận dụng bộ đệm và khó khôi phục cây từ dữ liệu
đệm sau sự cố như mất điện hoặc chương trình sụp đổ.

Ngoài ra, thư viện \texttt{\textless ext/rope\textgreater} của g++ đã có một
mẫu danh sách khối cơ bản: \texttt{\_\_gnu\_cxx::rope<T, Alloc>}. Vì vậy lập
trình viên C++ có thể dùng ngay một cấu trúc có sẵn.

\section*{Lời cảm ơn}

Tác giả cảm ơn thầy Liu Rujia đã góp ý về đề tài và nội dung bài viết; cô Hou
Xiaojing đã góp ý chỉnh sửa; Chen Yanhui đã cung cấp mã cho \textit{NOI 2005
Sequence}, giúp tác giả học thêm kinh nghiệm viết danh sách khối; Gao Yihan
đã giải thích cách ước lượng độ phức tạp; và Zhou Mengyu đã giới thiệu mẫu
danh sách khối trong thư viện \texttt{\textless ext/rope\textgreater}.

\section*{Tài liệu tham khảo và ghi chú}

\begin{itemize}
  \item Bài viết năm 2005 của Long Fan về ứng dụng của thứ tự.
  \item Bài viết năm 2007 của Yu Jiangwei về xử lý tốt các bài thống kê động.
  \item Lời giải \textit{NOI 2007} của Yu Linyun.
  \item Dữ liệu và chương trình chuẩn của \textit{Key Insertion} tại
  \url{http://neerc.ifmo.ru/past/2003.html}.
  \item Dữ liệu và chương trình chuẩn của \textit{Sort} tại
  \url{http://contest.felk.cvut.cz/07cerc/}.
  \item Giới thiệu về Vim của Bram Moolenaar tại
  \url{http://www.free-soft.org/FSM/english/issue01/vim.html}.
  \item Tài liệu về \texttt{rope} tại \url{http://www.sgi.com/tech/stl/} và
  \url{http://www.sgi.com/tech/stl/Rope.html}.
\end{itemize}

Một số ghi chú nhỏ trong bản gốc: còn có những cách cài đặt nhanh hơn, nhưng
cách mô phỏng bằng danh sách khối giúp giảm bớt nhánh xử lý và ít lỗi hơn;
bản dịch đề \textit{Sequence} lấy từ bài trại đông năm 2005 của Long Fan;
máy thử nghiệm của tác giả dùng CPU \(1.8\) GHz hai nhân và \(1\) GB RAM; mã
thử nghiệm nằm trong phần phụ lục của tài liệu gốc. Tác giả cũng ghi rằng có
thể dùng kỹ thuật rời rạc hóa thời gian thực cho một số dữ liệu của
\textit{Editor}; một phiên bản danh sách khối phức tạp hơn của tác giả qua
\(8\) test \textit{Necklace} và hơi quá thời gian ở test thứ \(9\); với
\textit{Happybirthday} của NOI 2006, mã danh sách khối ngắn hơn một chút so
với mã SBT của tác giả và qua được \(8\) test. Nếu dùng splay, có thể giảm độ
phức tạp mô phỏng xuống \(\Oh(\log n)\), đánh đổi bằng bộ nhớ lớn hơn.

\end{document}
